\chapter{Estado da Arte}

A previsão de séries temporais busca estimar valores futuros de um processo a partir de suas observações históricas ordenadas no tempo. Formalmente, a previsão direta $n$ passos à frente utiliza um vetor $\mathbf{X}_t$ de valores históricos da própria série e de eventuais variáveis exógenas para estimar valores futuros $Y_{t+n}$, buscando minimizar o erro preditivo esperado. Embora modelos estatísticos clássicos ofereçam boa tratabilidade analítica, sua dependência de premissas rígidas de estacionariedade faz com que frequentemente percam eficácia em séries com comportamentos não lineares e não estacionários, o que incentiva o uso de aprendizado de máquina na área \cite{masiniMachineLearningAdvances2021}.  

\textit{Ensemble learning} consolidou-se como abordagem fundamental em aprendizado de máquina ao combinar múltiplos modelos para superar o desempenho preditivo e a capacidade de genrealização de preditores individuais, mitango o risco de \textit{overfitting} \cite{fanDiverseModelsUnited2025}.
Conforme sitematizado por \citeonline{mendes-moreiraEnsembleApproachesRegression2012}, no contexto de regressão, a construção de um \textit{ensemble} é composta por três etapas: geração, onde se busca produzir um conjunto de modelos variados por meio de estratégias como \textit{Bagging}, \textit{Boosting} ou \textit{Stacking}; a seleção busca filtrar os preditores mais relevante e, por fim, a combinação, que agrega as previsões individuais por meio de diferentes métodos, como média ponderada, média aritmética ou meta-modelos. Ambas as revisões reforçam que a eficiência do \textit{ensemble} depende de maneira crítica da diversidade, apesar da promoção desta é apoiada em sua maior parte em estatégias heurísticas durante a geração e combinação \cite{fanDiverseModelsUnited2025, mendes-moreiraEnsembleApproachesRegression2012}

Na inteseção entre \textit{ensemble learning} e previsão de séries temporais, sistemas híbridos têm se destacado ao combinar modelos estatísticos e \textit{ensembles} para capturar a complexidade estrutural das séries temporais. Nesse sentido, \citeonline{deo.santosjuniorIntelligentHybridizationARIMA2019} propõem uma arquitetura hibrida em que um modelo linear (ARIMA) captura a dinâmica linear da série, enquanto um comitê modela os resíduos não lineares, os quais costumam apresentar heterocedasticidade e flutuações complexas. Nessa abordagem, a diversidade no \textit{ensemble} é promovida de forma implícita por meio de estratégias de subamostragem por meio de instâncias (\textit{Bagging}) e de atributos (\textit{Rnadom Subspace}) como também, são empregadas de maneira combinada (\textit{Random Patches}). Os autores demonstram ganhos empíricos sobre modelos individuais e híbridos tradicionais, apesar disso, a diversidade nesse sistema ainda é explorada como um artifício intuitivo de construção de agregações, a compreensão formal de como e sob quais condições a diversidade efetivamente contribui para a redução do erro preditivo permanece uma questão aberta na literatura.

Na tentativa de formalizar quantitativamente o papel da diversidade além da regressão quadrática, \citeonline{didaciDiversityClassifierEnsembles2013} derivaram decomposições do erro para \textit{ensembles} em classificação multiclasse sob perda 0/1, demonstrando que o desacordo relevante nesse cenário não reside nos rótulos escolhidos mas sim na correção das decisões individuais em relação ao \textit{ensemble}. Com essa formulação, os autores estendem os conceitos de "boa" e "má" diversidade: aumentar a divergência entre os modelos reduz o erro apenas nas amostras em que o ensemble acerta, mas penaliza o desempenho quando o \textit{ensemble} erra. O estudo também evidencia que as métricas tradicionais são apenas aproximações do termo exato de diversidade e não estabelecem um \textit{trade-off} claro com o erro médio individual. Apesar disso, conclui-se que o termo exato do estudo possui um carater essencialmente descritivo, pois, para seu cálculo, é necessário o conhecimento prévio do rótulo verdadeiro de $y$, o que já permitiria estimar a acurácia do \textit{ensemble}, por outro lado, a decomposição de ambiguidade na regressão sob perda quadrática prescinde de $y$, o que permite mensurar o desacordo mesmo em amostras não rotuladas.


Trazendo uma nova perspectiva para além da diversidade como uma heurística sem definição consensual, com exceção da perda quadrática com a combinação do \textit{ensemble} sendo a média aritmética, \citeonline{woodUnifiedTheoryDiversity} propõem uma teoria unificada derivada de primeiros princípios ao decompor diretamente funções de perda. Os autores revelam que a diversidade se trata de uma dimensão oculta na decomposição de viés-variância, formalizando ela como um componente para o ajuste do modelo que por sua vez quantifica as dependências estatísticas dos membros do \textit{ensemble}. Os autores trazem uma regra de agregação à função de perda por meio de um combinador centroide e provam decomposições aditivas em que o risco esperado equivale à soma do ruído, viés médio e variância média subtraída da diversidade, a qual quantifica o desacordo entre os preditores independente do rótulo real. Além disso, o estudo expande os achados de \cite{didaciDiversityClassifierEnsembles2013} ao generalizar a distinção entre "boa" e "má" diversidade para qualquer regra de combinação sob perdas não aditivas, comprovando que nesses casos, o efeito da diversidade é estritamente descritivo e dependente do rótulo real.

